\section{Results}
\label{sec:results}

\subsection{Real-transition training outperforms the imagination pipeline}
\label{sec:results-ppo}

With MLP torsos, real-transition training on the frozen JEPA encoders passes $0.94$ walls per episode
across nine runs, against $0.27$ across the $27$ imagined-training runs
(Figure~\ref{fig:contrast}; Appendix~\ref{app:round2}). This $0.67$-wall deficit belongs to the tested
imagination-and-training pipeline: dynamics, fitted heads and PPO protocol change together.
None of the imagined policies discovers the strategy; their maximum final survival is $3.1\%$.
The real frozen-latent policies also do not discover it, and their wall counts do not establish that
they use wall information.

True-pixel MLP policies pass $0.91$ walls per episode in three runs without discovery. This control
prevents attributing the frozen-latent MLP outcome to missing encoded information; it establishes
neither a representational ceiling nor a capacity ceiling for the MLP class. In the flow reference,
eight of nine policies pass $0.72$--$1.18$ walls and one passes $3.95$.
All share one dynamics checkpoint, so policy replication does not establish family-level superiority.
Training the gate JEPA for twice and four times the steps, $17{,}000$ and $34{,}000$, leaves in-model
PPO at $0$--$3\%$ survival (model seed $0$, policy seed $0$, each predictor variant).

\begin{figure}[t]
\centering
\renewcommand{\figureheight}{1.50in}
\paperfig{contrast.pdf}{1.0}
\caption{\textbf{The primary control contrast.} Final walls passed per real episode, with one dot per
policy and a horizontal mean marker. All groups use MLP torsos: imagined JEPA ($27$ policies),
real frozen latent ($9$), real pixels ($3$), and imagined flow ($9$, one checkpoint).
Marker shape identifies the JEPA encoder seed; the JEPA group pools its three predictor variants (original, residual sample, residual mean of eight) and the flow group its two sampling modes.
The real-transition advantage measures the complete pipeline contrast, including the PPO protocol.}
\label{fig:contrast}
\end{figure}

\subsection{Encoded information can support discovery}
\label{sec:results-wall}

A CNN trained in the real gate on decoded frozen latents discovers the strategy in $2/9$ runs,
both with encoder $1$ (Figure~\ref{fig:discovery}). The corresponding true-pixel CNN discovers it in
$2/3$ runs. Decoder fitting adds reconstruction supervision outside the world-model budget.
The result establishes useful information in at least one encoder under this read-out and learning
recipe, without asserting that every encoder preserves all task-relevant information.

Decoded real encodings retain the wall, but the original predicted rollouts lose it across all
three encoder seeds: by the second predicted step, decoded wall intensity is at most $0.02$
(Appendix~\ref{app:round2}). Figure~\ref{fig:hero} shows the contrast with a flow sample that retains
the wall and commits to a gap. A pixelwise mixture of the two gap futures would preserve shared
wall pixels; the observation differs from that explanation, but does not exclude averaging effects
in the learned latent model.

The fitted heads are a named competing explanation: on encoder $0$, termination AUC is $0.96$
for border deaths and $0.65$ for wall deaths, while predicted rewards are $-0.89$ and $-0.05$,
respectively, against the true $-1$; these are the signals on which the in-model policy trains.

\begin{figure}[t]
\centering
\renewcommand{\figureheight}{1.40in}
\paperfig{discovery.pdf}{1.0}
\caption{\textbf{Useful information remains encoded.} Final survival for true-pixel CNN policies and
CNN policies acting on decoded frozen latents, with marker shape identifying the encoder seed. Each dot is a policy run;
the dashed line marks the discovery criterion, survival above $50\%$. The decoded policies discover
the strategy in $2/9$ runs, both on encoder $1$, versus $2/3$ true-pixel runs. Decoder supervision is
off-budget. These are final endpoints, not training curves.}
\label{fig:discovery}
\end{figure}

\subsection{Scale and sampling bound the interpretation}
\label{sec:results-mean}

Counterfactual wall distance is smaller than total teacher-forced prediction error
(Figure~\ref{fig:geometry}). On $64$ pre-reveal contexts for encoder $0$, the first-step squared
distance between wall-present and wall-absent encodings is $0.035$, against a one-step residual
of $0.145$. The distance measures sensitivity to wall removal, whereas the residual includes all
predicted features. Their ordering is descriptive evidence, not a test that small amplitude causes
wall deletion or control failure. Raw distances across independently learned latent spaces are not
comparable information scales.

Sampling the fitted residual predictor does not rescue control: one sample and the mean of eight
each give $0.26$ walls per episode, against $0.27$ for the original predictor, with no discoveries.
The residual fitting protocol is described in Appendix~\ref{app:hparams}.
These tests neither implicate nor acquit the mean. They test a fitted distribution with shared
heads and training protocol, not an exact conditional mean or every distributional predictor.

\begin{figure}[t]
\centering
\renewcommand{\figureheight}{1.40in}
\paperfig{geometry.pdf}{1.0}
\caption{\textbf{A descriptive scale comparison.} Squared latent wall distance and total
teacher-forced one-step residual versus future step, averaged over the $64$ pre-reveal contexts
for JEPA model seed $0$. Wall distance compares real encoded futures with and without the wall,
under matched actions. Its position below total prediction error does not identify a cause of
wall deletion. No cross-model scale comparison is implied.}
\label{fig:geometry}
\end{figure}
